% GENERATED FILE. Edit canonical lessons, then run npm run generate:books.
% canonical-source-hash: fnv1a64:3bbdebd4d1534812
% canonical-lessons: FA-C284-mekanik, FA-C284-arayeshgar, FA-C284-postchi, FA-C284-ketabdar, FA-C284-hesabdar

\chapter{Mechanic, Hairdresser, Postman, Librarian, Accountant}
\label{ch:fa-mechanic-hairdresser-postman-librarian-accountant}
\hlchaptermodality{284}

\begin{chapteropening}
\textbf{By the end of this chapter:} \emph{I can say a mechanic, a hairdresser, a postman, a librarian and an accountant.}

Complete the last lesson of chapter 284: I can say a mechanic, a hairdresser, a postman, a librarian and an accountant.
\end{chapteropening}

\section[\texorpdfstring{\fa{مکانیک} (mekânik)}{مکانیک (mekânik)}]{\fa{مکانیک} (mekânik) — a mechanic}

\label{lesson:FA-C284-mekanik}



\begin{quote}
Before the new one: say the Persian for a screwdriver, then the Persian for pliers.
\end{quote}

\subsection*{You'll want to know: \fa{مکانیک}}
\textbf{\fa{مکانیک}} — \emph{mekânik} — \textquotedblleft{}a mechanic\textquotedblright{}.

\begin{cousinweb}
\emph{mekânik} is a loan, usually traced to French \emph{mécanique}, \textquotedblleft{}mechanics\textquotedblright{}.
\end{cousinweb}

\begin{grammarlens}[title={What you've built}]
One more word for this chapter.
\end{grammarlens}

\subsection*{Your turn}
\begin{itemize}
\raggedright
  \item \emph{Say it:} \emph{mekânik}
  \item \emph{Say it:} \emph{mekânik}, once more
  \item \emph{Say it:} say \emph{mekânik}
  \item \emph{Recall:} say the Persian for a screwdriver, then the Persian for pliers, then say \emph{mekânik} again
\end{itemize}

\begin{tcolorbox}[breakable,colback=teal!4,colframe=teal!35!black,title={Before you move on}]
What does \fa{مکانیک} mean? (\textquotedblleft{}A mechanic\textquotedblright{}.) Say it once more.
\end{tcolorbox}
\section[\texorpdfstring{\fa{آرایشگر} (ârâyeshgar)}{آرایشگر (ârâyeshgar)}]{\fa{آرایشگر} (ârâyeshgar) — a hairdresser}

\label{lesson:FA-C284-arayeshgar}



\begin{quote}
Before the new one: say the Persian for pliers, then the Persian for a mechanic.
\end{quote}

\subsection*{You'll want to know: \fa{آرایشگر}}
\textbf{\fa{آرایشگر}} — \emph{ârâyeshgar} — \textquotedblleft{}a hairdresser\textquotedblright{}.

\begin{cousinweb}
\emph{ârâyeshgar} is built from \emph{ârâyesh}, \textquotedblleft{}adornment\textquotedblright{} (from \emph{ârâstan}, \textquotedblleft{}to adorn\textquotedblright{}), and \emph{-gar}, \textquotedblleft{}doer\textquotedblright{}: the adorner. Both are native Persian, inherited rather than borrowed.
\end{cousinweb}

\begin{grammarlens}[title={What you've built}]
One more word for this chapter.
\end{grammarlens}

\subsection*{Your turn}
\begin{itemize}
\raggedright
  \item \emph{Say it:} \emph{ârâyeshgar}
  \item \emph{Say it:} \emph{ârâyeshgar}, once more
  \item \emph{Say it:} say \emph{ârâyeshgar}
  \item \emph{Recall:} say the Persian for pliers, then the Persian for a mechanic, then say \emph{ârâyeshgar} again
\end{itemize}

\begin{tcolorbox}[breakable,colback=teal!4,colframe=teal!35!black,title={Before you move on}]
What does \fa{آرایشگر} mean? (\textquotedblleft{}A hairdresser\textquotedblright{}.) Say it once more.
\end{tcolorbox}
\section[\texorpdfstring{\fa{پستچی} (postchi)}{پستچی (postchi)}]{\fa{پستچی} (postchi) — a postman}

\label{lesson:FA-C284-postchi}



\begin{quote}
Before the new one: say the Persian for a mechanic, then the Persian for a hairdresser.
\end{quote}

\subsection*{You'll want to know: \fa{پستچی}}
\textbf{\fa{پستچی}} — \emph{postchi} — \textquotedblleft{}a postman\textquotedblright{}.

\begin{cousinweb}
\emph{postchi} is \emph{post}, \textquotedblleft{}mail\textquotedblright{}, borrowed from French \emph{poste}, with the agent ending \emph{-chi}, which Persian took from Turkic: a French word with a Turkic tail.
\end{cousinweb}

\begin{grammarlens}[title={What you've built}]
One more word for this chapter.
\end{grammarlens}

\subsection*{Your turn}
\begin{itemize}
\raggedright
  \item \emph{Say it:} \emph{postchi}
  \item \emph{Say it:} \emph{postchi}, once more
  \item \emph{Say it:} say \emph{postchi}
  \item \emph{Recall:} say the Persian for a mechanic, then the Persian for a hairdresser, then say \emph{postchi} again
\end{itemize}

\begin{tcolorbox}[breakable,colback=teal!4,colframe=teal!35!black,title={Before you move on}]
What does \fa{پستچی} mean? (\textquotedblleft{}A postman\textquotedblright{}.) Say it once more.
\end{tcolorbox}
\section[\texorpdfstring{\fa{کتابدار} (ketâbdâr)}{کتابدار (ketâbdâr)}]{\fa{کتابدار} (ketâbdâr) — a librarian}

\label{lesson:FA-C284-ketabdar}



\begin{quote}
Before the new one: say the Persian for a hairdresser, then the Persian for a postman.
\end{quote}

\subsection*{You'll want to know: \fa{کتابدار}}
\textbf{\fa{کتابدار}} — \emph{ketâbdâr} — \textquotedblleft{}a librarian\textquotedblright{}.

\begin{cousinweb}
\emph{ketâbdâr} is \emph{ketâb}, \textquotedblleft{}book\textquotedblright{}, with \emph{-dâr}, \textquotedblleft{}holder\textquotedblright{}: the book-keeper. \emph{ketâb} is borrowed from Arabic \emph{kitāb} (root \emph{k-t-b}, Lane); \emph{-dâr} is native Persian.
\end{cousinweb}

\begin{grammarlens}[title={What you've built}]
One more word for this chapter.
\end{grammarlens}

\subsection*{Your turn}
\begin{itemize}
\raggedright
  \item \emph{Say it:} \emph{ketâbdâr}
  \item \emph{Say it:} \emph{ketâbdâr}, once more
  \item \emph{Say it:} say \emph{ketâbdâr}
  \item \emph{Recall:} say the Persian for a hairdresser, then the Persian for a postman, then say \emph{ketâbdâr} again
\end{itemize}

\begin{tcolorbox}[breakable,colback=teal!4,colframe=teal!35!black,title={Before you move on}]
What does \fa{کتابدار} mean? (\textquotedblleft{}A librarian\textquotedblright{}.) Say it once more.
\end{tcolorbox}
\section[\texorpdfstring{\fa{حسابدار} (hesâbdâr)}{حسابدار (hesâbdâr)}]{\fa{حسابدار} (hesâbdâr) — an accountant}

\label{lesson:FA-C284-hesabdar}



\begin{quote}
Before the new one: say the Persian for a postman, then the Persian for a librarian.
\end{quote}

\subsection*{You'll want to know: \fa{حسابدار}}
\textbf{\fa{حسابدار}} — \emph{hesâbdâr} — \textquotedblleft{}an accountant\textquotedblright{}.

\begin{cousinweb}
\emph{hesâbdâr} is \emph{hesâb}, \textquotedblleft{}account\textquotedblright{}, with \emph{-dâr}, \textquotedblleft{}holder\textquotedblright{}: the account-keeper. \emph{hesâb} is borrowed from Arabic \emph{ḥisāb}, \textquotedblleft{}reckoning\textquotedblright{} (root \emph{ḥ-s-b}, Lane); \emph{-dâr} is native.
\end{cousinweb}

\begin{grammarlens}[title={What you've built}]
One more word for this chapter.
\end{grammarlens}

\subsection*{Your turn}
\begin{itemize}
\raggedright
  \item \emph{Say it:} \emph{hesâbdâr}
  \item \emph{Say it:} \emph{hesâbdâr}, once more
  \item \emph{Say it:} say \emph{hesâbdâr}
  \item \emph{Recall:} say the Persian for a postman, then the Persian for a librarian, then say \emph{hesâbdâr} again
\end{itemize}

\begin{tcolorbox}[breakable,colback=teal!4,colframe=teal!35!black,title={Before you move on}]
What does \fa{حسابدار} mean? (\textquotedblleft{}An accountant\textquotedblright{}.) Say it once more.
\end{tcolorbox}
